%======================================================================
% CAPÍTULO 12 - ARQUITECTURA DEL SISTEMA
%======================================================================
\chapter{Arquitectura del Sistema}

\bigskip

\dropcap{L}{a} arquitectura de Sentinel se fundamenta en principios de diseño que priorizan el rendimiento, la mantenibilidad y la trazabilidad. Este capítulo detalla los componentes del sistema y sus interacciones.

\bigskip

\section{Visión General de la Arquitectura}

Sentinel implementa una Arquitectura Hexagonal con las siguientes capas:

\begin{enumerate}
    \item Capa de Interfaz: CLI y Manifiestos
    \item Capa de Orquestación: Kernel
    \item Capa de Negocio: Motor de Cálculo y Fusión
    \item Capa de Datos: gRPC Client y Serialización
    \item Capa de Presentación: Exportación CSV
\end{enumerate}

\bigskip

\section{Componentes Principales}

\subsection{CLI (Interfaz de Línea de Comandos)}

El punto de entrada del sistema. Responsabilidades:

\begin{itemize}
    \item Parseo de argumentos de línea de comandos.
    \item Carga y validación de manifiestos JSON.
    \item Inicialización del Kernel.
    \item Presentación de resultados y estadísticas.
\end{itemize}

\bigskip

Ejemplo de ejecución:

\begin{lstlisting}
cargo run -- -x -m nomina_2026.json
\end{lstlisting}

\bigskip

\subsection{Kernel (Orquestador)}

El centro de procesamiento que coordina todas las operaciones:

\begin{enumerate}
    \item Gestión del ciclo de ejecución.
    \item Orquestación de la carga de datos.
    \item Coordinación de cálculos.
    \item Control de la exportación.
    \item Telemetría y logging.
\end{enumerate}

\bigskip

\subsection{Cargador (gRPC Client)}

Módulo responsable de:

\begin{itemize}
    \item Conexión gRPC al servidor Sandra.
    \item Streaming de datos desde la base de datos.
    \item Deserialización paralela de JSON/NDJSON.
    \item Normalización de datos a estructuras Rust.
\end{itemize}

\bigskip

El cargamento de datos se realiza en paralelo mediante spawn de tareas asíncronas, permitiendo que mientras se procesa un lote, el siguiente ya está siendo descargado.

\bigskip

\subsection{Motor Rhai}

Motor de scripting para evaluación de fórmulas de primas:

\begin{itemize}
    \item Compilación de expresiones Rhai desde la base de datos.
    \item Ejecución en sandbox aislado.
    \item Integración con Rayon para paralelismo.
    \item Manejo robusto de errores.
\end{itemize}

\bigskip

\subsection{Motor de Cálculo (Calculadora Nativa)}

Implementa los algoritmos de cálculo de nómina en Rust nativo:

\begin{itemize}
    \item Cálculo de antigüedad y tiempo de servicio.
    \item Determinación de sueldo base.
    \item Cálculo de alícuotas.
    \item Asignaciones y garantías.
\end{itemize}

\vfill
\newpage
